% Reviewed translation: PDF pages 271--276 / printed pages 257--262.
% The original scan is authoritative; mathematical tokens were checked against it.
\MathBookSourcePage{271}
\section{三维向量势--涡量格式}
\label{sec:incompressible-vector-potential-vorticity}

把前面各节的格式推广到三维 Stokes 问题并不容易.
显然的原因是, 确定向量势的条件比确定二维流函数的条件复杂得多.
因此这里只把 Section~\ref{sec:incompressible-sfvp}
的 stream function--vorticity 格式推广到
$\mathbb R^3$ 中一个很简单区域上的齐次 Stokes 问题.

本节始终假设 $\Omega$ 是 $\mathbb R^3$ 中有界单连通开集,
其边界 $\Gamma$ 是连通、有界的多面体.
压力逼近留到最后一小节处理. 首先降低与双调和问题有关的函数空间正则性;
使用 $H(\operatorname{curl};\Omega)$ 中的函数已经足够.
这将导出 $H(\operatorname{curl};\Omega)$ 中协调有限元的有限维子空间,
它们不是 $H^1(\Omega)^3$ 的子空间.
最后, 因为使用间断元, 压力也采用与
Subsection~\ref{subsec:incompressible-hhj-discontinuous-pressure}
非常相似的间断逼近.

\subsection{三维 Stokes 问题的混合表述}
\label{subsec:incompressible-vector-potential-mixed-formulation}

给定 $\mathbf f\in L^2(\Omega)^3$, 考虑齐次 Stokes 问题:
求 $(\mathbf u,p)\in H_0^1(\Omega)^3\times L_0^2(\Omega)$, 使
\begin{equation}\label{eq:incompressible-vector-potential-stokes}
\left\{
\begin{aligned}
-\nu\Delta\mathbf u+\operatorname{grad}p&=\mathbf f,\\
\operatorname{div}\mathbf u&=0
\end{aligned}
\right\}\quad\text{于 }\Omega,
\qquad
\mathbf u=0\quad\text{于 }\Gamma.
\end{equation}
该问题可解释为速度向量势 $\boldsymbol\psi$ 的双调和问题,
$\mathbf u=\operatorname{curl}\boldsymbol\psi$, 且
\begin{equation}\label{eq:incompressible-vector-potential-space}
\Psi=\{\boldsymbol\phi\in L^2(\Omega)^3;
\ \operatorname{div}\boldsymbol\phi\in H_0^1(\Omega),
\ \operatorname{curl}\boldsymbol\phi\in H_0^1(\Omega)^3,
\ \boldsymbol\phi\times\mathbf n|_\Gamma=0\}.
\end{equation}

为导出混合表述, 用 $\operatorname{curl}\boldsymbol\phi$
检验 \eqref{eq:incompressible-vector-potential-stokes}.

\MathBookSourcePage{272}
因为 $\operatorname{div}\mathbf u=0$, 首先得到
\[
\nu\langle\operatorname{curl}\operatorname{curl}\mathbf u,
\operatorname{curl}\boldsymbol\phi\rangle
=(\mathbf f,\operatorname{curl}\boldsymbol\phi)
\qquad\forall\boldsymbol\phi\in H_0(\operatorname{curl};\Omega),
\]
其中括号表示 $H_0(\operatorname{curl};\Omega)$ 与其对偶之间的对偶配对.
令
\begin{equation}\label{eq:incompressible-vector-potential-vorticity}
\boldsymbol\omega=\operatorname{curl}\mathbf u.
\end{equation}
若 $\boldsymbol\omega\in H(\operatorname{curl};\Omega)$, 则
\[
\nu(\operatorname{curl}\boldsymbol\omega,
\operatorname{curl}\boldsymbol\phi)
=(\mathbf f,\operatorname{curl}\boldsymbol\phi)
\quad\forall\boldsymbol\phi\in H_0(\operatorname{curl};\Omega).
\]
又因 $\mathbf u=\operatorname{curl}\boldsymbol\psi$,
$\operatorname{div}\boldsymbol\psi=0$, $\boldsymbol\psi\in\Psi$,
\eqref{eq:incompressible-vector-potential-vorticity} 等价于
\[
(\operatorname{curl}\operatorname{curl}\boldsymbol\psi,\boldsymbol\mu)
=(\boldsymbol\omega,\boldsymbol\mu)
\qquad\forall\boldsymbol\mu\in L^2(\Omega)^3.
\]
把 $\boldsymbol\mu$ 限制到 $H(\operatorname{curl};\Omega)$, 得
\[
(\operatorname{curl}\boldsymbol\psi,
\operatorname{curl}\boldsymbol\mu)
=(\boldsymbol\omega,\boldsymbol\mu)
\qquad\forall\boldsymbol\mu\in H(\operatorname{curl};\Omega).
\]
因此问题 $(\widetilde Q)$ 为: 求
$(\boldsymbol\psi,\boldsymbol\omega)
\in H_0(\operatorname{curl};\Omega)\times
H(\operatorname{curl};\Omega)$, 使
\begin{align}
\nu(\operatorname{curl}\boldsymbol\omega,
\operatorname{curl}\boldsymbol\phi)
&=(\mathbf f,\operatorname{curl}\boldsymbol\phi)
&&\forall\boldsymbol\phi\in H_0(\operatorname{curl};\Omega),
\label{eq:incompressible-vector-potential-mixed-first}\\
(\operatorname{curl}\boldsymbol\psi,
\operatorname{curl}\boldsymbol\mu)
&=(\boldsymbol\omega,\boldsymbol\mu)
&&\forall\boldsymbol\mu\in H(\operatorname{curl};\Omega),
\label{eq:incompressible-vector-potential-mixed-second}\\
\operatorname{div}\boldsymbol\psi&=0.
\label{eq:incompressible-vector-potential-divergence-free}
\end{align}
它至少有解 $(\boldsymbol\psi,\boldsymbol\omega=-\Delta\boldsymbol\psi)$.
反之, 若齐次第一式成立, 在第二式中依次取
$\boldsymbol\phi=\boldsymbol\psi$ 和
$\boldsymbol\mu=\boldsymbol\omega$, 得
$\boldsymbol\omega=0$ 和
$\operatorname{curl}\boldsymbol\psi=0$.
再由 $\operatorname{div}\boldsymbol\psi=0$ 及 $\Gamma$ 连通,
可知 $\boldsymbol\psi=0$.

\begin{Theorem}\label{thm:incompressible-vector-potential-continuous-uniqueness}
若问题 \eqref{eq:incompressible-vector-potential-stokes} 的解满足
$\operatorname{curl}\mathbf u\in H(\operatorname{curl};\Omega)$,
则 \eqref{eq:incompressible-vector-potential-mixed-first}--
\eqref{eq:incompressible-vector-potential-divergence-free}
有唯一解
$(\boldsymbol\psi,-\Delta\boldsymbol\psi)$,
其中 $\operatorname{curl}\boldsymbol\psi=\mathbf u$.
\end{Theorem}

令
\[
\Phi_0=\{\boldsymbol\phi\in H_0(\operatorname{curl};\Omega);
\ \operatorname{div}\boldsymbol\phi=0\}.
\]
问题 $(\widetilde Q)$ 等价于求
$(\boldsymbol\psi,\boldsymbol\omega)
\in\Phi_0\times H(\operatorname{curl};\Omega)$, 使
\begin{equation}\label{eq:incompressible-vector-potential-reduced-problem}
\left\{
\begin{aligned}
\nu(\operatorname{curl}\boldsymbol\omega,
\operatorname{curl}\boldsymbol\phi)
&=(\mathbf f,\operatorname{curl}\boldsymbol\phi)
&&\forall\boldsymbol\phi\in\Phi_0,\\
(\operatorname{curl}\boldsymbol\psi,
\operatorname{curl}\boldsymbol\mu)
&=(\boldsymbol\omega,\boldsymbol\mu)
&&\forall\boldsymbol\mu\in H(\operatorname{curl};\Omega).
\end{aligned}
\right.
\end{equation}

\MathBookSourcePage{273}
引入
\[
\widetilde X=\{\operatorname{curl}\boldsymbol\phi;
\ \boldsymbol\phi\in\Phi_0\}\times L^2(\Omega)^3,
\qquad \widetilde M=H(\operatorname{curl};\Omega),
\]
以及, 对
$u=(\operatorname{curl}\boldsymbol\psi,\boldsymbol\omega)$,
$v=(\operatorname{curl}\boldsymbol\phi,\boldsymbol\theta)$,
\[
\widetilde a(u,v)=\nu(\boldsymbol\omega,\boldsymbol\theta),
\quad
\widetilde b(v,\boldsymbol\mu)
=(\operatorname{curl}\boldsymbol\phi,
\operatorname{curl}\boldsymbol\mu)
-(\boldsymbol\theta,\boldsymbol\mu),
\quad
\langle l,v\rangle=(\mathbf f,\operatorname{curl}\boldsymbol\phi).
\]
则问题写成
\begin{equation}\label{eq:incompressible-vector-potential-abstract-problem}
\left\{
\begin{aligned}
\widetilde a(u,v)+\widetilde b(v,\lambda)&=\langle l,v\rangle
&&\forall v\in\widetilde X,\\
\widetilde b(u,\boldsymbol\mu)&=0
&&\forall\boldsymbol\mu\in\widetilde M.
\end{aligned}
\right.
\end{equation}
相应核空间为
\[
\widetilde V=\{v=(\operatorname{curl}\boldsymbol\phi,\boldsymbol\theta)
\in\widetilde X;
(\operatorname{curl}\boldsymbol\phi,
\operatorname{curl}\boldsymbol\mu)
=(\boldsymbol\theta,\boldsymbol\mu)
\ \forall\boldsymbol\mu\in H(\operatorname{curl};\Omega)\}.
\]
由
\begin{equation}\label{eq:incompressible-vector-potential-poincare}
\|\boldsymbol\phi\|_{0,\Omega}
\leq C\|\operatorname{curl}\boldsymbol\phi\|_{0,\Omega}
\qquad\forall\boldsymbol\phi\in\Phi_0,
\end{equation}
可在 $\widetilde X$ 上取范数
\[
\|v\|_{\widetilde X}
=\bigl(\|\operatorname{curl}\boldsymbol\phi\|_{0,\Omega}^2
+\|\boldsymbol\theta\|_{0,\Omega}^2\bigr)^{1/2}.
\]
在 $\widetilde V$ 上还有
\begin{equation}\label{eq:incompressible-vector-potential-kernel-control}
\|\operatorname{curl}\boldsymbol\phi\|_{0,\Omega}
\leq C\|\boldsymbol\theta\|_{0,\Omega}.
\end{equation}
所以 $|v|=\|\boldsymbol\theta\|_{0,\Omega}$
是 $\widetilde V$ 上与 $\widetilde X$ 范数等价的范数, 且
\[
\widetilde a(v,v)=\nu|v|^2
\geq\frac{\nu}{C^2+1}\|v\|_{\widetilde X}^2.
\]

\MathBookSourcePage{274}
同样容易验证弱 inf--sup 条件
\[
\sup_{v\in\widetilde X}
\frac{\widetilde b(v,\boldsymbol\mu)}{\|v\|_{\widetilde X}}
\geq\|\boldsymbol\mu\|_{0,\Omega}
\qquad\forall\boldsymbol\mu\in\widetilde M.
\]
这里 Lagrange 乘子满足 $\lambda=\nu\boldsymbol\omega$.

\begin{Remark}\label{rem:incompressible-vector-potential-two-dimensional-analogy}
问题 \eqref{eq:incompressible-vector-potential-reduced-problem}
与二维 stream function--vorticity 问题相似.
\end{Remark}

\begin{Remark}\label{rem:incompressible-vector-potential-larger-test-space}
\eqref{eq:incompressible-vector-potential-abstract-problem} 的第一式
实际上在比 $\widetilde X$ 更大的空间上成立, 即
$v=(\operatorname{curl}\boldsymbol\phi,\boldsymbol\theta)$,
$(\boldsymbol\phi,\boldsymbol\theta)
\in H_0(\operatorname{curl};\Omega)\times L^2(\Omega)^3$.
\end{Remark}

\subsection{$H(\operatorname{curl};\Omega)$ 中的混合逼近}
\label{subsec:incompressible-vector-potential-mixed-approximation}

引入有限维空间
\begin{equation}\label{eq:incompressible-vector-potential-discrete-spaces}
\Phi_h\subset H_0(\operatorname{curl};\Omega),
\qquad M_h\subset H(\operatorname{curl};\Omega),
\qquad\Theta_h\subset H_0^1(\Omega),
\end{equation}
并假设 $\Phi_h\subset M_h$.
由于 $\Phi_h$ 未必包含于 $H(\operatorname{div};\Omega)$,
无散条件写成
\begin{equation}\label{eq:incompressible-vector-potential-discrete-divergence-free}
(\boldsymbol\phi_h,\operatorname{grad}q_h)=0
\qquad\forall q_h\in\Theta_h.
\end{equation}
即用
\begin{equation}\label{eq:incompressible-vector-potential-discrete-kernel}
\Phi_{h0}=\{\boldsymbol\phi_h\in\Phi_h;
\ \boldsymbol\phi_h\text{ 满足 }
\eqref{eq:incompressible-vector-potential-discrete-divergence-free}\}
\end{equation}
逼近 $\Phi_0$. 要求存在 $C^*>0$, 使
\begin{equation}\label{eq:incompressible-vector-potential-discrete-poincare}
\|\boldsymbol\phi_h\|_{0,\Omega}
\leq C^*\|\operatorname{curl}\boldsymbol\phi_h\|_{0,\Omega}
\qquad\forall\boldsymbol\phi_h\in\Phi_{h0}.
\end{equation}
离散问题 $(Q_h)$ 为求
$(\boldsymbol\psi_h,\boldsymbol\omega_h)
\in\Phi_{h0}\times M_h$, 使
\begin{equation}\label{eq:incompressible-vector-potential-constrained-discrete-problem}
\left\{
\begin{aligned}
\nu(\operatorname{curl}\boldsymbol\omega_h,
\operatorname{curl}\boldsymbol\phi_h)
&=(\mathbf f,\operatorname{curl}\boldsymbol\phi_h)
&&\forall\boldsymbol\phi_h\in\Phi_{h0},\\
(\operatorname{curl}\boldsymbol\psi_h,
\operatorname{curl}\boldsymbol\mu_h)
&=(\boldsymbol\omega_h,\boldsymbol\mu_h)
&&\forall\boldsymbol\mu_h\in M_h.
\end{aligned}
\right.
\end{equation}

\MathBookSourcePage{275}
由于对检验函数有约束, 实际计算改用整个 $\Phi_h$:
\begin{equation}\label{eq:incompressible-vector-potential-unconstrained-discrete-problem}
\left\{
\begin{aligned}
\nu(\operatorname{curl}\boldsymbol\omega_h,
\operatorname{curl}\boldsymbol\phi_h)
&=(\mathbf f,\operatorname{curl}\boldsymbol\phi_h)
&&\forall\boldsymbol\phi_h\in\Phi_h,\\
(\operatorname{curl}\boldsymbol\psi_h,
\operatorname{curl}\boldsymbol\mu_h)
&=(\boldsymbol\omega_h,\boldsymbol\mu_h)
&&\forall\boldsymbol\mu_h\in M_h.
\end{aligned}
\right.
\end{equation}
两问题等价需要附加假设: 对每个 $\boldsymbol\phi_h\in\Phi_h$,
存在 $\boldsymbol\phi_{h0}\in\Phi_{h0}$ 使
\begin{equation}\label{eq:incompressible-vector-potential-curl-lifting}
\operatorname{curl}\boldsymbol\phi_h
=\operatorname{curl}\boldsymbol\phi_{h0}.
\end{equation}

\begin{Lemma}\label{lem:incompressible-vector-potential-discrete-solvability}
令 \eqref{eq:incompressible-vector-potential-discrete-spaces} 成立,
$\Phi_h\subset M_h$, 并以
\eqref{eq:incompressible-vector-potential-discrete-kernel}
定义 $\Phi_{h0}$. 在
\eqref{eq:incompressible-vector-potential-discrete-poincare} 下,
问题 $(Q_h)$ 在 $\Phi_{h0}\times M_h$ 中有唯一解.
若此外 \eqref{eq:incompressible-vector-potential-curl-lifting} 成立,
则 \eqref{eq:incompressible-vector-potential-unconstrained-discrete-problem}
与 $(Q_h)$ 等价.
\end{Lemma}

为应用抽象误差定理, 取
\begin{align}
X_h&=\{\operatorname{curl}\boldsymbol\phi_h;
\ \boldsymbol\phi_h\in\Phi_{h0}\}\times M_h,
\label{eq:incompressible-vector-potential-discrete-product-space}\\
V_h&=\{(\operatorname{curl}\boldsymbol\phi_h,\boldsymbol\theta_h)
\in X_h;
(\operatorname{curl}\boldsymbol\phi_h,
\operatorname{curl}\boldsymbol\mu_h)
=(\boldsymbol\theta_h,\boldsymbol\mu_h)
\ \forall\boldsymbol\mu_h\in M_h\}.
\label{eq:incompressible-vector-potential-discrete-nullspace}
\end{align}
虽然 $V_h$ 一般不包含于 $\widetilde V$, 但由
\eqref{eq:incompressible-vector-potential-discrete-poincare},
\begin{equation}\label{eq:incompressible-vector-potential-discrete-kernel-control}
\|\operatorname{curl}\boldsymbol\phi_h\|_{0,\Omega}
\leq C^*\|\boldsymbol\theta_h\|_{0,\Omega}
\quad\forall(\operatorname{curl}\boldsymbol\phi_h,
\boldsymbol\theta_h)\in V_h.
\end{equation}
故 $|v_h|=\|\boldsymbol\theta_h\|_{0,\Omega}$
是 $V_h$ 上的等价范数, 可直接应用抽象定理.

\MathBookSourcePage{276}
\begin{Theorem}\label{thm:incompressible-vector-potential-abstract-error}
保留 Lemma~\ref{lem:incompressible-vector-potential-discrete-solvability}
除 \eqref{eq:incompressible-vector-potential-curl-lifting} 外的全部假设.
则 $(Q_h)$ 的解满足
\begin{equation}\label{eq:incompressible-vector-potential-vorticity-abstract-error}
\begin{split}
\|\boldsymbol\omega-\boldsymbol\omega_h\|_{0,\Omega}
\leq{}&2\inf_{v_h\in V_h}|u-v_h|\\
&+(1+C^{*2})^{1/2}
\inf_{\boldsymbol\mu_h\in M_h}
\|\boldsymbol\omega-\boldsymbol\mu_h\|_{H(\operatorname{curl};\Omega)},
\end{split}
\end{equation}
以及
\begin{equation}\label{eq:incompressible-vector-potential-velocity-abstract-error}
\begin{split}
\|\operatorname{curl}(\boldsymbol\psi-\boldsymbol\psi_h)\|_{0,\Omega}
\leq{}&(1+C^{*2})^{1/2}
\left\{\inf_{v_h\in V_h}\|u-v_h\|_{\widetilde X}\right.\\
&\left.\qquad+C^*\inf_{\boldsymbol\mu_h\in M_h}
\|\boldsymbol\omega-\boldsymbol\mu_h\|_{H(\operatorname{curl};\Omega)}\right\},
\end{split}
\end{equation}
其中 $C^*$ 是
\eqref{eq:incompressible-vector-potential-discrete-poincare} 的常数.
\end{Theorem}

现在需要估计 $V_h$ 的逼近误差. 二维情形的抽象论证仍然成立.
\begin{Lemma}\label{lem:incompressible-vector-potential-kernel-approximation}
在 Lemma~\ref{lem:incompressible-vector-potential-discrete-solvability}
的记号下, 对每个
$v=(\operatorname{curl}\boldsymbol\phi,\boldsymbol\theta)\in\widetilde V$,
\begin{equation}\label{eq:incompressible-vector-potential-kernel-approximation}
\inf_{w_h\in V_h}|v-w_h|
\leq\inf_{v_h=(\operatorname{curl}\boldsymbol\phi_h,
\boldsymbol\theta_h)\in X_h}
\left\{2|v-v_h|+
\sup_{\boldsymbol\mu_h\in M_h}
\frac{(\operatorname{curl}(\boldsymbol\phi-\boldsymbol\phi_h),
\operatorname{curl}\boldsymbol\mu_h)}
{\|\boldsymbol\mu_h\|_{0,\Omega}}\right\}.
\end{equation}
把右端的半范数换成 $\widetilde X$ 范数,
同样可得到 $\|v-w_h\|_{\widetilde X}$ 的上界.
\end{Lemma}
